\subsection{Lifting through coefficient ideals with exchange upper rings}
We use ordinary exchange theory with its established scope.
The fixed-cardinality exchange property passes to finite direct
sums and direct summands, and a module has the two-summand
exchange property exactly when its endomorphism ring satisfies
the exchange equation. Complete proofs are given in
\cite[Lemmas~2.3--2.4, Proposition~2.7 and Corollary~2.8,
pp.~2--6]{NielsenExchange}. The ideal-lifting implication is
\cite[Proposition~1.1 and Corollary~1.3, p.~270]{NicholsonExchange}.
We include the short calculation needed here.

The terminology also has a global semiring convention.
Das and Kar require absorbing complements for every
multiplicative idempotent in their definition of an exchange
semiring \cite[Section~2 and Definition~3.11, pp.~4 and~10]{DasKarExchange}.
Our hypothesis is the ordinary exchange equation in each
ring $F(e)$; no identification of these conditions is used.

\begin{proof}[Proof of Theorem~\ref{thm:exchange-coefficient-lifting}]
First, if a unital ring $B$ has a finite complete orthogonal
family $p_i$ and every $p_iBp_i$ is exchange, then $B$ is
exchange. Indeed, the right modules $p_iB$ have endomorphism
rings $p_iBp_i$, by left multiplication. The cited module
criterion and finite-sum property show that
$B_B=\bigoplus_i p_iB$ has the two-summand exchange property.
The same criterion applies to its endomorphism ring $B$.

In an exchange ring, let $b^2-b$ belong to a two-sided ideal
$M$. Choose $p^2=p=rb$ and $1-p=s(1-b)$ by the exchange
equation. Then
\[
 p-b=p(1-b)-(1-p)b=(r-s)(b-b^2)\in M.
\]
Thus idempotents lift modulo every two-sided ideal. The
calculation also works for left ideals, though that extra scope
is not needed below.

Fix $E^2=E$ and put $D=I\vee E$. Then
\[
 D^2=D,\qquad DE=E=ED.
\]
Write $A=A_E(T)$ and $B=A_D(T)$. The latter is a unital ring
whose coordinate idempotents have diagonal corners $F(d_{ii})$.
Each $d_{ii}=1\vee e_{ii}$ is an upper multiplicative
idempotent, since $e_{ii}^2\leq e_{ii}$. By the hypothesis
and the preceding finite-corner argument, $B$ is exchange.

The upward transition gives a ring map $\phi:A\to B$.
Its kernel $K$ annihilates $A$ on both sides. For example,
if $\phi(x)=0$ and $y\in A$, transition compatibility and
$DE=E$ identify $xy$ with the mixed product $\phi(x)y$ in
the fiber at $E$, hence with zero. The other order uses
$ED=E$. These are ordinary group zeros at their specified
supports; no vanishing assumption on $F(0)$ is used.

Put $N=\prod_{i,j}\mathcal I_{e_{ij}}\subseteq A$.
Componentwise quotient identifies $A/N$ with
$A_E(T/\mathcal I)$. The image $M=\phi(N)$ is a two-sided
ideal of $B$. Indeed, a mixed product of a $D$-supported
coefficient matrix and an $E$-supported matrix in $N$ again
has support $E$ and coefficients in $\mathcal I$, by
$DE=ED=E$ and compatibility of the coefficient ideal.
Applying $\phi$ proves both ideal closures. This argument
requires no injectivity of the transitions.

Let $\bar a\in A/N$ be idempotent and choose a lift
$a\in A$. Then $\phi(a)$ is idempotent modulo $M$.
The ordinary exchange calculation gives $p^2=p\in B$
with $p-\phi(a)\in M$. Choose $n\in N$ with
$\phi(n)=p-\phi(a)$ and set $a_0=a+n$. Its defect
$a_0^2-a_0$ lies in $K$, so the annihilator property gives
\[
 a_0^3=a_0^2,\qquad (a_0^2)^2=a_0^2.
\]
The idempotent $a_0^2\in A_E(T)$ reduces to
$\bar a^2=\bar a$. All operations stay in that exact
additive matrix fiber, including entries whose coefficients
are group zeros. This proves the projector-lifting assertion.

A finite projective over $T/\mathcal I$ is the split image
of such a projector. Its lift has a finite-projective split
image over $T$, and extension of scalars preserves this split
retract of a finite free semimodule. Reduction respects direct
sums, so it surjects on $V$ and on its group completion.
For $\mathcal I=C(T)$, the unconditional isomorphism-detection
part of Theorem~\ref{thm:upper-jacobson-quotient} also gives
injectivity on $V$. Its bijection is a monoid isomorphism and
therefore induces the asserted $K_0$-isomorphism.
\end{proof}

The ordinary exchange property is used for the whole finite
upper coefficient-matrix ring. Scalar idempotent lifting in
each diagonal ring alone would not justify that step. The
theorem concerns existence of projective lifts; it does not
supply the local decompositions or uniqueness used in the
semiperfect amplitude classification.
